\section{Logikgatter — die Bausteine des Prozessors}
\label{sec:grundlagen:logikgatter}

Bevor wir verstehen, wie Computer rechnen, müssen wir wissen, woraus sie gebaut sind.
Das Herzstück eines Computers ist der Prozessor (die CPU). Er besteht aus vielen
kleinen Schaltungen, die mit nur zwei Zuständen arbeiten: \textbf{0} und \textbf{1}.

Diese Schaltungen nennt man \textbf{Logikgatter}. Die wichtigsten sind:

\begin{description}
  \item[UND-Gatter (AND)]\hfill\\
    Das UND-Gatter liefert nur dann eine 1, wenn beide Eingänge 1 sind.
    \begin{center}
    \begin{tabular}{c|c|c}
    A & B & A AND B \\
    \hline
    0 & 0 & 0 \\
    0 & 1 & 0 \\
    1 & 0 & 0 \\
    1 & 1 & 1 \\
    \end{tabular}
    \end{center}

  \item[ODER-Gatter (OR)]\hfill\\
    Das ODER-Gatter liefert eine 1, wenn mindestens einer der Eingänge 1 ist.
    \begin{center}
    \begin{tabular}{c|c|c}
    A & B & A OR B \\
    \hline
    0 & 0 & 0 \\
    0 & 1 & 1 \\
    1 & 0 & 1 \\
    1 & 1 & 1 \\
    \end{tabular}
    \end{center}

  \item[XOR-Gatter (exklusives ODER)]\hfill\\
    Das XOR-Gatter liefert eine 1, wenn genau einer der Eingänge 1 ist.
    \begin{center}
    \begin{tabular}{c|c|c}
    A & B & A XOR B \\
    \hline
    0 & 0 & 0 \\
    0 & 1 & 1 \\
    1 & 0 & 1 \\
    1 & 1 & 0 \\
    \end{tabular}
    \end{center}

  \item[NICHT-Gatter (NOT)]\hfill\\
    Das NOT-Gatter invertiert den Eingang: aus 0 wird 1, aus 1 wird 0.
    \begin{center}
    \begin{tabular}{c|c}
    A & NOT A \\
    \hline
    0 & 1 \\
    1 & 0 \\
    \end{tabular}
    \end{center}
\end{description}

Aus diesen einfachen Gattern kann man alles bauen – sogar eine Addition!
% \begin{exercise}
% ...
% \end{exercise}